% steps/step09.tex
\section[گام نهم: بهبود و بهینه‌سازی]{\faTools\ \ گام نهم: بهبود و بهینه‌سازی}
\begin{stepbox}

\field{جلوگیری از بیش‌برازش و حفظ‌کردن داده‌ها (\lr{Regularization})}{
    شبکه‌های عصبی عمیق با پارامترهای بالا مستعد بیش‌برازش (\lr{Overfitting}) هستند؛ به این معنا که به‌جای یادگیری الگوها و ویژگی‌های کلی تومور، داده‌های آموزشی را حفظ می‌کنند. جهت مقابله با این پدیده، از تکنیک‌های منظم‌سازی (\lr{Regularization}) مانند \lr{Dropout} استفاده می‌شود. در این روش، با خاموش کردن تصادفی درصدی از نرون‌ها در هر مرحله از آموزش، شبکه ناچار می‌شود مسیرهای موازی و ویژگی‌های قوی‌تری را برای تصمیم‌گیری پیدا کند و تعمیم‌پذیری آن به شکل چشمگیری افزایش یابد.
}

\field{تکنیک افزایش داده در زمان تست (\lr{TTA})}{
    روش \lr{Test-Time Augmentation (TTA)} برای افزایش پایداری و دقت پیش‌بینی مدل بر روی داده‌های آزمون اعمال می‌شود. در این تکنیک، هر تصویر ورودی جدید علاوه بر حالت اصلی، تحت چند تبدیل ساده و هندسی (مانند قرینه‌سازی افقی یا چرخش‌های جزئی) قرار گرفته و به مدل تغذیه می‌شود. در نهایت، پیش‌بینی نهایی از میانگین‌گیری یا تلفیق خروجی‌های به‌دست‌آمده از این زوایای مختلف حاصل می‌شود که این امر خطاهای موضعی مدل را کاهش می‌دهد.
}

\field{استفاده از خرد جمعی مدل‌ها (\lr{Ensemble Learning})}{
    به منظور بهینه‌سازی نهایی و افزایش ضریب اطمینان بالینی، می‌توان از یادگیری جمعی (\lr{Ensemble}) بهره برد. در این رویکرد، چندین معماری یا مدل مختلف که با استراتژی‌ها و بخش‌های متفاوتی از داده‌ها آموزش دیده‌اند، با یکدیگر ترکیب می‌شوند. در زمان استنتاج، خروجی تمام مدل‌ها برای یک تصویر مشخص محاسبه شده و با استفاده از روش‌هایی نظیر میانگین‌گیری وزن‌دار یا رای‌گیری پیکسلی، دقیق‌ترین ماسک نهایی تولید می‌شود که این کار ریسک خطای تک‌مدلی را به حداقل می‌رساند.
}

\end{stepbox}
